% ─────────────────────────────────────────────────────────────────────────────
\section{The Gap Between Competence and Understanding}
\label{sec:intro_gap}
% ─────────────────────────────────────────────────────────────────────────────

The past decade of machine learning research has been defined by a striking
asymmetry. Progress on benchmarks has been rapid and measurable. Progress on
the underlying question of what learned representations actually encode has
been slow and contested. The two trajectories have now diverged to a degree
that demands explanation.

Contemporary deep learning systems achieve results that would have seemed
implausible a decade ago. Large language models pass professional licensing
examinations, synthesise code from natural-language specifications, and
produce coherent text across domains that require substantial world
knowledge~\cite{brown_gpt3, wei_emergent}. Vision systems trained on
internet-scale data match or exceed human performance on curated image
classification benchmarks~\cite{dosovitskiy_vit, zhai_scaling}. Reinforcement
learning agents master games of enormous strategic depth, from Go to
StarCraft~\cite{silver_alphago, vinyals_starcraft}. By the standard measures
of the field, progress has been exceptional.

Yet when these same systems are placed outside the conditions under which they
were evaluated, a different picture emerges. An agent that reliably completes
multi-step tasks on familiar software interfaces fails when the interface
changes, when the task sequence is novel, or when an unexpected state requires
a response not present in its training distribution~\cite{acu_survey}. A
language model that solves competition-level mathematics cannot reliably apply
a simple spatial rule to a grid of objects it has not seen
before~\cite{cogitao}. A world model that accurately predicts one step ahead
diverges catastrophically over a ten-step horizon, accumulating errors that
render its predictions useless for planning~\cite{koopman_jepa}. These are not
isolated failures attributable to insufficient data or model capacity. They
form a pattern.

The pattern is this: current deep learning systems exhibit high competence
within the statistical regime of their training distribution and brittle
behaviour outside it. They are sensitive to surface features in ways that
suggest their internal representations do not encode the structural properties
of their environment. They fail at tasks that require composing learned
concepts in configurations not encountered during training. And they degrade
under distribution shift in ways that are disproportionate to the magnitude
of the shift. These observations are not new, and individual instances of them
have been documented extensively in the
literature~\cite{geirhos_shortcut, lake_generalization, marcus_deep_learning}.
What has been lacking is a unifying account of why they occur and what, if
anything, can be done about them.

This thesis provides such an account. We argue that the failures listed above
are not incidental consequences of imperfect training or insufficient scale.
They are structural consequences of what gradient-based optimisation on a
fixed objective actually produces: representations whose geometry is
determined by the statistical structure of the training signal rather than
by the causal structure of the world. This claim is not merely
programmatic. In Chapter~\ref{chap:selfattention} we establish it
mathematically for the case of self-attention, the representational mechanism
at the core of modern Transformer architectures, and validate it empirically
across multiple model families and input modalities. From this diagnosis,
two distinct lines of response follow, each explored in depth in the
chapters that follow.


% ─────────────────────────────────────────────────────────────────────────────
\section{The Standard Explanation and Why It Is Insufficient}
\label{sec:intro_scaling}
% ─────────────────────────────────────────────────────────────────────────────

The dominant response to the observed failure modes has been to appeal to
scale. The empirical scaling laws of~\cite{kaplan_scaling, hoffmann_chinchilla}
show that model performance on standard benchmarks improves predictably with
increases in parameter count, dataset size, and compute budget. Under this
view, the failures described above are temporary artefacts of systems that are
not yet large enough: with sufficient scale, generalisation will follow.

We take this argument seriously. The empirical evidence for scaling is
substantial, and dismissing it would be intellectually dishonest. The question
is not whether scale produces improvements, but whether it addresses the
specific class of failures that motivates this thesis. We argue that it does
not, and we present two bodies of evidence in support of this claim.

The first concerns the relationship between model capacity and practical
usability. The compression methods currently required to deploy large models in
resource-constrained settings, including quantisation and low-rank
adaptation~\cite{dettmers_qlora, hu_lora}, do not merely reduce a model's
memory footprint. They degrade its reasoning capabilities in ways that are
disproportionate to the degree of compression applied~\cite{private_llm}. In a
systematic benchmark covering seven methods across two major families of
efficient inference technique, we find that quantised models exhibit
qualitative deterioration in text quality and structured reasoning that is
detectable both by automated metrics and by human evaluators, even at
quantisation levels that are routinely described as lossless in the
literature~\cite{private_llm}. This finding matters for the scaling argument
because it reveals that the capabilities produced by large-scale training are
not robustly encoded: they are sensitive to the precise arithmetic used to
represent the weights. A representation that encodes genuine structural
knowledge of its domain should not be so fragile.

The second body of evidence concerns compositional generalisation. If the
failures observed in deployed systems were primarily a scale problem, one would
expect them to diminish as model capacity increases, and to be
particularly severe in architectures that lack specific structural inductive
biases. We find the opposite. Using the COGITAO framework~\cite{cogitao},
a controlled evaluation environment for compositional and systematic
generalisation in object-centric visual reasoning, we test architectures that
incorporate precisely the inductive biases that compositionality is thought to
require: diffusion-based Transformers and recurrent architectures with
adaptive computation time. Despite strong in-domain performance, these models
fail systematically when task rules are composed in configurations not seen
during training~\cite{cogitao}. The failure is not ameliorated by the
inductive biases; it appears to be a deeper property of how gradient-based
learning encodes multi-step relational structure.

Taken together, these findings suggest that the failures are not primarily
a quantitative shortfall in scale or a qualitative shortfall in architecture
design. They point instead to something more fundamental: a mismatch between
what gradient-based optimisation on a fixed objective produces and what is
required for reliable performance on tasks that demand structured understanding.


% ─────────────────────────────────────────────────────────────────────────────
\section{A Structural Diagnosis: Training Objectives Shape Representations}
\label{sec:intro_diagnosis}
% ─────────────────────────────────────────────────────────────────────────────

To understand the source of these failures, it is necessary to examine what
backpropagation-trained systems actually learn. The standard account describes
learning as function approximation: given a dataset and a loss function, gradient
descent finds parameters that minimise the loss, and the resulting function
generalises to new inputs to the extent that those inputs are drawn from the
same distribution as the training data. This account is accurate as far as it
goes. The question it leaves open is what structure the learned function has
beyond its loss value, and whether that structure is appropriate for the tasks
to which the system will be applied.

We address this question directly in Chapter~\ref{chap:selfattention}, where we
develop a mathematical framework for analysing the internal structure of
self-attention matrices in trained Transformer models. Self-attention is not
merely a component of modern architectures; it is the primary mechanism by
which these models organise and retrieve information, and understanding its
representational geometry is essential for understanding what these models
know and how they know it.

Our central finding is the following. The structural properties of learned
self-attention weight matrices are determined by the training objective, not
by the content of the training data. Bidirectional training objectives, of the
kind used in encoder models such as BERT and ModernBERT, induce symmetry in the
self-attention weight matrices. Autoregressive training objectives, of the kind
used in decoder models such as GPT, LLaMA, and Mistral, induce directionality
and column dominance. These structural signatures are not noise; they are
mathematically derivable from the gradient flow induced by each objective, and
they are empirically consistent across model families and across input
modalities including text, vision, and audio~\cite{selfattention_structures}.

The implication of this result for the question of generalisation is direct.
If the geometry of a model's representations is determined by the structure
of its training objective, then the representations encode the statistical
patterns that the objective is sensitive to, not the causal or structural
properties of the domain from which the training data was drawn. A model
trained with a next-token prediction objective does not, as a consequence of
that training, learn that objects in the world have stable identities, that
causes precede effects, or that rules compose in a specific algebraic
structure. It learns to predict the next token. The representations it develops
are optimised for that task and are well-calibrated within the statistical
regime where that task is well-defined. Outside that regime, they carry no
guarantees.

This diagnosis is not a claim that Transformer models are useless or that they
learn nothing of value. They clearly do learn a great deal. The claim is more
specific: the structure of what they learn is shaped by the objective, and
the failure modes observed empirically are predictable consequences of this
shaping. Compositional generalisation fails because the training objective does
not make the model sensitive to compositional structure; it makes the model
sensitive to co-occurrence statistics. Long-horizon planning fails because no
standard training objective rewards the model for maintaining internally
consistent predictions over extended sequences of interactions. Distributional
robustness fails because the representations are calibrated to the training
distribution, not to the underlying generative process that produced it.

This diagnosis raises a precise research question: given that these failures
are structural properties of what gradient-based learning on standard objectives
produces, what can be done? We pursue two distinct lines of response in this
thesis, which we introduce in the following section.


% ─────────────────────────────────────────────────────────────────────────────
\section{Two Lines of Response}
\label{sec:intro_responses}
% ─────────────────────────────────────────────────────────────────────────────

The diagnosis of Section~\ref{sec:intro_diagnosis} admits two logically
distinct responses. The first accepts the basic framework of gradient-based
learning on fixed objectives and asks whether hard mathematical constraints,
imposed architecturally or through domain knowledge, can correct specific
failure modes before they manifest. The second asks whether the representational
substrate must itself change: whether the class of architectures that
gradient-based learning is applied to needs to be fundamentally reconceived.
We pursue both, and we are explicit about what each can and cannot achieve.

\subsection{Hard Constraints as Structural Corrections}
\label{subsec:intro_constraints}

If a failure mode can be identified precisely enough to be formalised, it can
often be corrected by imposing a constraint that prevents it. This is the
principle underlying the first line of response, and we instantiate it
across three domains in Chapters~\ref{chap:koopman}
and~\ref{chap:constraints}.

The most precise instance concerns multi-step prediction in learned world
models. In any nonlinear world model, approximation errors compound across
rollout steps: each step applies a function with some Lipschitz constant
greater than one, and errors at each step are amplified by that constant in
subsequent steps, producing exponential divergence over long
horizons~\cite{koopman_jepa}. This failure is not a consequence of
insufficient model capacity or imperfect training. It is a mathematical
property of nonlinear function composition. It cannot be corrected by scaling
or by improved optimisation, because it is structural.

We address this failure in Chapter~\ref{chap:koopman} by introducing
Koopman-JEPA, a self-supervised world model that lifts complex visual
observations into a latent space where the dynamics of the environment are
represented by a globally linear operator. In this lifted space, multi-step
prediction reduces to iterated matrix multiplication, and error accumulation
is bounded to grow linearly with horizon rather than exponentially. We develop
three variants with progressively stronger spectral constraints on the
transition matrix, culminating in a Cayley parameterisation that
architecturally enforces strict stability and satisfies the stabilisability
precondition of discrete-time linear quadratic regulation. The empirical
results are correspondingly precise: Koopman variants reduce long-horizon
prediction error by up to 99\% on action-free tasks and 91\% on
action-conditioned tasks at a ten-step horizon, compared to a baseline that
diverges catastrophically at the same horizon~\cite{koopman_jepa}.

The constraint principle is not limited to dynamics. In Chapter~\ref{chap:constraints},
we examine two further cases where structural knowledge, encoded as a hard
constraint, compensates for what unconstrained gradient-based learning fails
to recover. In the domain of clinical image segmentation, encoding
tissue-specific Hounsfield unit ranges as binary input masks, a constraint
derived directly from the biophysics of X-ray attenuation documented in the
medical literature, improves segmentation accuracy by up to 5\% on
intramuscular adipose tissue and enables a model trained on 50\% of the
available data to match the performance of one trained on the
full dataset~\cite{hu_ranges}. In the domain of information retrieval, a
hybrid pipeline that combines lexical precision, dense semantic retrieval, and
large-language-model-based re-ranking, effectively grounding inference in a
structured external knowledge source at test time, achieves first place on the
CLEF CheckThat! 2025 development leaderboard and third place on the hidden
test set without any external training data~\cite{deep_retrieval}.

These three cases illustrate different forms of the same principle. In each,
a specific, identifiable failure mode of unconstrained gradient-based learning
is corrected by encoding structural knowledge that the training objective is
not sensitive to. The constraint operates at a different level from the
objective: it does not change what the model is trained to do, but it changes
the space of solutions available to it.

We are explicit about the limits of this approach. Constraint engineering
requires that the structure of the failure mode be known in advance, that it
be formalised precisely enough to be encoded architecturally or in the input
representation, and that the constraint be compatible with gradient-based
optimisation. For many of the most consequential failure modes of current AI
systems, including generalisation to novel compositional structures and
reliable behaviour under large distributional shift, these conditions are
difficult to satisfy. Correcting these failures requires not a better
constraint but a different kind of representation.

\subsection{Representational Reconception}
\label{subsec:intro_paradigm}

The second line of response begins from the observation that the failure modes
are not merely failures of a particular architecture or training procedure.
They are failures of the representational principle common to virtually all
current deep learning: the encoding of information in dense, distributed
activation patterns over unstructured weight matrices, updated globally by
gradient descent. Systems of this kind can represent anything, in the sense
that they are universal function approximators. But representational universality
is not the same as representational appropriateness. A representation that
can encode any function provides no guarantee that the function it has learned
is the right one, in the sense of capturing the stable, compositional, invariant
structure that tasks outside the training distribution require.

Biological neural systems take a different approach. The visual system, to take
the most studied case, does not represent visual scenes as dense, unstructured
activation patterns over all its neurons simultaneously. It organises
representations hierarchically and compositionally, with neurons that respond
to local features assembled into larger structural units that respond to
progressively more abstract properties of the
scene~\cite{hubel_receptive, felleman_hierarchy}. The representations are
sparse, structured, and invariant to a range of transformations by construction,
not as a consequence of training on a sufficiently large dataset with a
sufficiently expressive architecture.

In Chapter~\ref{chap:cna} we introduce the Cooperative Network Architecture
(CNA), which proposes a fundamental departure from distributed representation
toward structured assembly. Rather than encoding sensory input as a pattern of
activation over a fixed weight matrix, CNA represents its input through a
dynamically assembled network of neurons, termed a \textit{net}, constructed
from overlapping \textit{net fragments} learned from statistical regularities
in the input without supervision. Because nets are compositional by
construction, novel inputs are represented by assembling fragments in new
configurations rather than by interpolating within a learned embedding space.
This representational principle produces properties that are absent from
standard deep learning architectures: figure completion, noise resilience, and
generalisation to out-of-distribution inputs emerge from the structure of the
representation rather than from loss engineering~\cite{cna}.

CNA does not inherit the scaling infrastructure of standard deep learning.
The dynamic assembly of structured networks is not straightforwardly
expressible as dense matrix multiplication, and it does not currently benefit
from the hardware and software optimisations that make large-scale gradient
descent practical. This is a genuine limitation. But it is important not to
conflate the scalability of an approach with its correctness. The evidence
presented in this thesis suggests that the approach currently at scale is
not producing the representational properties required for robust generalisation.
CNA demonstrates that those properties can be produced by a different
representational principle, establishing a proof of concept for a class of
architectures worth developing further.

The relationship between the two lines of response is not competitive.
Constraint engineering within the current paradigm and representational
reconception beyond it operate at different timescales and address different
aspects of the problem. The constraint-based approaches of
Chapters~\ref{chap:koopman} and~\ref{chap:constraints} offer immediate,
deployable corrections to specific and well-understood failure modes.
The architectural approach of Chapter~\ref{chap:cna} addresses the underlying
representational question that constraint engineering leaves open.
Both are necessary parts of a serious response to the diagnosis of
Section~\ref{sec:intro_diagnosis}.


% ─────────────────────────────────────────────────────────────────────────────
\section{Why This Matters}
\label{sec:intro_stakes}
% ─────────────────────────────────────────────────────────────────────────────

The failures described in this thesis are not purely of theoretical interest.
As AI systems are increasingly deployed in settings where they must act
autonomously over extended time horizons, on tasks that were not anticipated
at training time, and in environments that differ from their training conditions,
the gap between benchmark competence and reliable behaviour becomes consequential.

Our survey of agents for computer use~\cite{acu_survey}, covering 87 research
papers and 33 benchmark datasets, documents this gap directly. The systems
that currently achieve the strongest results on ACU benchmarks are evaluated
under conditions that differ systematically from those of real-world deployment:
tasks are shorter, interfaces are more uniform, and failure cases are less
diverse. When deployment conditions are examined, the gap between benchmark
performance and practical utility is substantial. The six research gaps we
identify, including insufficient generalisation, limited planning horizons,
and a disconnect between research evaluation and deployment constraints,
are not incidental properties of immature technology. They are predictable
consequences of the structural limitations diagnosed in this thesis.

The stakes extend beyond the computer use setting. In clinical AI, where
segmentation models are trained on datasets that are expensive to acquire and
annotate, the inability to generalise from limited data under distributional
shift has direct consequences for patient outcomes~\cite{hu_ranges}. In
scientific literature retrieval, where the task requires bridging informal
natural-language descriptions and formal academic knowledge bases, the failure
of purely parametric models to ground their outputs in precise, verifiable
sources limits their utility for high-stakes
applications~\cite{deep_retrieval}. The argument of this thesis, that these
failures share a structural cause and admit structured responses, is therefore
not only of scientific interest. It has implications for the conditions under
which AI systems can be deployed responsibly.


% ─────────────────────────────────────────────────────────────────────────────
\section{Research Questions}
\label{sec:intro_questions}
% ─────────────────────────────────────────────────────────────────────────────

The preceding sections motivate four research questions that organise the
contributions of this thesis.

\begin{enumerate}

    \item \textbf{What are the concrete, systematic failure modes of
    state-of-the-art deep learning systems when deployed in demanding
    real-world settings?} This question motivates Part~\ref{part:failures},
    where we survey agents for computer use and examine two independent failure
    domains to establish an empirical foundation for the rest of the thesis.

    \item \textbf{What structural properties of learned representations
    underlie these failures?} This question motivates
    Part~\ref{part:diagnosis}, where we develop a mathematical framework
    for analysing self-attention geometry and show that the structure of learned
    representations is determined by the training objective.

    \item \textbf{Can hard inductive constraints, imposed architecturally or
    through domain knowledge, correct specific failure modes within the existing
    deep learning paradigm?} This question motivates
    Part~\ref{part:constraints}, where we introduce Koopman-JEPA and examine
    two complementary applications of domain-structural constraints.

    \item \textbf{Can neuro-inspired representational principles address the
    limitations that constraint engineering within the current paradigm cannot?}
    This question motivates Part~\ref{part:paradigm}, where we introduce the
    Cooperative Network Architecture and examine the representational properties
    it produces.

\end{enumerate}


% ─────────────────────────────────────────────────────────────────────────────
\section{Contributions of This Thesis}
\label{sec:intro_contributions}
% ─────────────────────────────────────────────────────────────────────────────

The primary contributions of this thesis are as follows.

\paragraph{A comprehensive survey and taxonomy of agents for computer use.}
We survey 87 research papers and 33 benchmark datasets on ACUs, introduce a
multifaceted taxonomy structured along domain, interaction, and agent
perspectives, and identify six persistent research gaps that characterise the
current state of the field~\cite{acu_survey}. Beyond its value as a reference
for the ACU community, this survey establishes the empirical basis for the
thesis argument: the identified gaps are symptomatic of structural properties
of the underlying learning paradigm.

\paragraph{A mathematical framework for analysing self-attention geometry.}
We derive the structures governing weight updates in self-attention matrices
under different training objectives, prove that bidirectional training induces
symmetry and autoregressive training induces directionality and column
dominance, and validate these theoretical results empirically across
ModernBERT, GPT, LLaMA3, and Mistral on text, vision, and audio
modalities~\cite{selfattention_structures}. We further show that symmetric
initialisation improves encoder performance on language tasks, demonstrating
that the theoretical analysis yields actionable insights.

\paragraph{Koopman-JEPA: a self-supervised world model with provably linear
error accumulation.}
We introduce Koopman-JEPA, which lifts visual observations into a latent space
governed by a globally linear operator, bounding multi-step prediction error
to grow linearly with horizon rather than exponentially~\cite{koopman_jepa}.
We develop three variants with progressively stronger spectral constraints,
culminating in a Cayley parameterisation that satisfies the stabilisability
precondition of discrete-time linear quadratic regulation. We provide both
theoretical bounds and extensive empirical validation, showing up to 99\% and
91\% reduction in prediction error on action-free and action-conditioned tasks
respectively, and approximately 16-fold improvement in sample efficiency at
10\% data availability.

\paragraph{Structured inductive biases for data-efficient learning and
inference-time grounding.}
We demonstrate that encoding tissue-specific Hounsfield unit ranges as binary
input masks improves CT segmentation accuracy and data efficiency in
resource-constrained clinical settings~\cite{hu_ranges}. We further demonstrate
that a hybrid retrieval pipeline combining lexical, semantic, and
cross-encoder re-ranking achieves state-of-the-art performance on scientific
literature retrieval from informal social media descriptions without external
training data~\cite{deep_retrieval}. Together with Koopman-JEPA, these
contributions establish that the constraint principle generalises across
domains and structural levels.

\paragraph{The Cooperative Network Architecture.}
We introduce CNA, which represents sensory input through dynamically assembled
networks of neurons built from compositional fragments learned without
supervision~\cite{cna}. CNA demonstrates that figure completion, noise
resilience, and out-of-distribution generalisation can be produced by a
representational principle rather than by loss engineering, establishing
a foundation for neuro-inspired architectures that embed structural invariance
by design.


% ─────────────────────────────────────────────────────────────────────────────
\section{Thesis Structure}
\label{sec:intro_structure}
% ─────────────────────────────────────────────────────────────────────────────

The remainder of this thesis is organised as follows.

Chapter~\ref{chap:background} provides the technical background required to
follow the subsequent chapters. It covers the foundations of gradient-based
learning and Transformer architectures, the theory of Koopman operators and
their relevance to dynamical systems learning, and the neuro-inspired concepts
that motivate the Cooperative Network Architecture.

\textbf{Part~\ref{part:failures}: Identifying Failures.}
Chapter~\ref{chap:acu} presents the comprehensive survey and taxonomy of agents
for computer use, documenting the six research gaps and arguing that they are
symptomatic of structural limitations in the current learning paradigm.
Chapter~\ref{chap:contextual} synthesises supporting evidence from two
independent domains: the benchmark of efficient large language models and the
evaluation of compositional generalisation using COGITAO.

\textbf{Part~\ref{part:diagnosis}: Diagnosing the Root Cause.}
Chapter~\ref{chap:selfattention} presents the mathematical framework for
self-attention geometry, establishes the theoretical results on symmetry and
directionality, and discusses the implications for understanding what
gradient-based learning on standard objectives produces.

\textbf{Part~\ref{part:constraints}: Mitigating Failures with Hard Constraints.}
Chapter~\ref{chap:koopman} introduces Koopman-JEPA, presenting both the
theoretical analysis of error accumulation in nonlinear world models and the
full empirical evaluation of the Koopman variants.
Chapter~\ref{chap:constraints} examines two complementary instantiations of
the constraint principle: Hounsfield unit masks for clinical segmentation and
hybrid retrieval for scientific literature grounding.

\textbf{Part~\ref{part:paradigm}: A Change of Paradigm.}
Chapter~\ref{chap:cna} introduces the Cooperative Network Architecture,
presents its representational principle, and evaluates the structural
properties it produces. Chapter~\ref{chap:dfc} presents the second
neuro-inspired contribution of this thesis~[to be completed].

Chapter~\ref{chap:discussion} synthesises the findings across all four parts,
discusses the trade-offs between constraint engineering and representational
reconception, and situates the thesis contributions within the broader
trajectory of AI research. Chapter~\ref{chap:conclusion} states the
conclusions of the thesis.

\bigskip

\noindent A note on authorship. This thesis is cumulative: it is based on
published and submitted research articles produced in collaboration with
co-authors. All primary contributions, those presented as full chapters, are
works in which the thesis author holds a first or shared-first authorship
position. Supporting contributions, presented as shorter sections in
Chapters~\ref{chap:contextual} and~\ref{chap:constraints}, are works with
more limited author involvement. In all cases, the specific contribution of
the thesis author is stated at the opening of the relevant chapter or section.
The framing, argumentation, and synthesis that connect the individual works
into the thesis narrative are the sole work of the thesis author.